\documentclass[11pt,a4paper]{article}

% ------------------------------------------------------------------------------
% Packages
% ------------------------------------------------------------------------------
\usepackage[utf8]{inputenc}
\usepackage[T1]{fontenc}
\usepackage[french]{babel}
\usepackage{lmodern}
\usepackage[a4paper,margin=2.2cm]{geometry}
\usepackage{xcolor}
\usepackage{booktabs}
\usepackage{array}
\usepackage{float}
\usepackage{caption}
\usepackage{listings}
\usepackage{tikz}
\usetikzlibrary{babel,arrows.meta}
\usepackage{hyperref}

\hypersetup{colorlinks=true, linkcolor=bleu, urlcolor=bleu}

% ------------------------------------------------------------------------------
% Couleurs
% ------------------------------------------------------------------------------
\definecolor{bleu}{HTML}{1E3A8A}
\definecolor{chemin}{HTML}{2563EB}
\definecolor{cheminfond}{HTML}{DBEAFE}
\definecolor{vert}{HTML}{16A34A}
\definecolor{vertfond}{HTML}{DCFCE7}
\definecolor{orange}{HTML}{EA580C}
\definecolor{orangefond}{HTML}{FFEDD5}
\definecolor{rouge}{HTML}{DC2626}
\definecolor{gris}{HTML}{64748B}
\definecolor{codefond}{HTML}{F8FAFC}
\definecolor{codecadre}{HTML}{CBD5E1}

% ------------------------------------------------------------------------------
% Style du code C++
% ------------------------------------------------------------------------------
\lstset{
  language=C++,
  basicstyle=\ttfamily\footnotesize,
  keywordstyle=\color{chemin}\bfseries,
  commentstyle=\color{gris}\itshape,
  stringstyle=\color{vert},
  backgroundcolor=\color{codefond},
  frame=single,
  rulecolor=\color{codecadre},
  breaklines=true,
  breakatwhitespace=true,
  columns=fullflexible,
  keepspaces=true,
  showstringspaces=false,
  tabsize=4,
  numbers=left,
  numberstyle=\tiny\color{gris},
  captionpos=b,
}
\renewcommand{\lstlistingname}{Code}

% ------------------------------------------------------------------------------
% Styles TikZ
% ------------------------------------------------------------------------------
\tikzset{
  noeud/.style={circle, draw=black!75, thick, fill=white, minimum size=8mm,
                inner sep=0pt, font=\small\bfseries},
  actif/.style={noeud, draw=orange, fill=orangefond},
  neuf/.style={noeud, draw=vert, fill=vertfond},
  arete/.style={draw=black!45, thick},
  echange/.style={<->, >={Stealth[length=2.5mm]}, rouge, very thick},
  case/.style={rectangle, draw=black!60, thick, minimum width=8mm, minimum height=7mm,
               inner sep=0pt, font=\small\bfseries},
  noeudcase/.style={fill=white},
  actifcase/.style={draw=orange, fill=orangefond},
  neufcase/.style={draw=vert, fill=vertfond},
  indice/.style={font=\scriptsize, text=gris},
}

\setlength{\parindent}{0pt}
\setlength{\parskip}{5pt}
\renewcommand{\arraystretch}{1.15}

\title{\textbf{Tas max}\\[4pt]
       \large Insertions et extractions sur 200\,000 valeurs, comparaison avec la STL}
\author{Zakariae Bakkari\\
        \small Programmation Objet Avancée et Design Patterns\\
        \small\url{https://github.com/zakariae-bakkari/cpp_ilisi2}}
\date{\today}

\begin{document}

\maketitle

% ==============================================================================
\section{Objectif}
% ==============================================================================

Implémenter un \textbf{tas max} en C++ (\texttt{src/10\_tas\_max.cpp}), y insérer les
200\,000 valeurs du fichier \texttt{data/data.txt}, puis extraire le maximum jusqu'à vider
le tas, en mesurant le temps et le nombre d'échanges de chaque phase. Puis comparer avec
trois structures de la STL : \texttt{std::priority\_queue}, \texttt{std::set} et
\texttt{std::map}.

\textbf{Code source :}
\begin{itemize}
  \item notre tas (option \textbf{11} du menu \texttt{./run.sh}, ou \texttt{make run\_10}) :\\
        \url{https://github.com/zakariae-bakkari/cpp_ilisi2/blob/main/iterators/src/10_tas_max.cpp}
  \item comparaison STL (option \textbf{12}, ou \texttt{make run\_11}) :\\
        \url{https://github.com/zakariae-bakkari/cpp_ilisi2/blob/main/iterators/src/11_tas_max_comparaison_stl.cpp}
\end{itemize}

% ==============================================================================
\section{Un arbre rangé dans un tableau}
% ==============================================================================

Un tas max est un arbre binaire \textbf{complet} (rempli niveau par niveau, de gauche à
droite) qui respecte une seule règle : \textbf{chaque parent est plus grand que ses
enfants}. Le maximum est donc toujours à la racine.

Comme l'arbre est complet, on le range dans un simple \texttt{vector<int>}, sans pointeurs.
Pour la case $i$ :

\begin{center}
\begin{tabular}{l l}
parent        & $(i - 1) / 2$ \\
enfant gauche & $2i + 1$ \\
enfant droit  & $2i + 2$ \\
\end{tabular}
\end{center}

\begin{figure}[H]
\centering
\begin{minipage}{0.5\textwidth}\centering
\begin{minipage}[t]{0.32\textwidth}\centering
\begin{tikzpicture}[scale=0.62, transform shape]
  \draw[arete] (0,0) -- (-1.6,-1.1);
  \draw[arete] (0,0) -- (1.6,-1.1);
  \draw[arete] (-1.6,-1.1) -- (-2.4,-2.2);
  \draw[arete] (-1.6,-1.1) -- (-0.8,-2.2);
  \draw[arete] (1.6,-1.1) -- (0.8,-2.2);
  \node[noeud] (n0) at (0,0) {30};
  \node[noeud] (n1) at (-1.6,-1.1) {15};
  \node[noeud] (n2) at (1.6,-1.1) {20};
  \node[noeud] (n3) at (-2.4,-2.2) {4};
  \node[noeud] (n4) at (-0.8,-2.2) {8};
  \node[noeud] (n5) at (0.8,-2.2) {10};
  \node[case, noeudcase] at (-2.00,-3.4) {30};
  \node[indice] at (-2.00,-3.95) {0};
  \node[case, noeudcase] at (-1.20,-3.4) {15};
  \node[indice] at (-1.20,-3.95) {1};
  \node[case, noeudcase] at (-0.40,-3.4) {20};
  \node[indice] at (-0.40,-3.95) {2};
  \node[case, noeudcase] at (0.40,-3.4) {4};
  \node[indice] at (0.40,-3.95) {3};
  \node[case, noeudcase] at (1.20,-3.4) {8};
  \node[indice] at (1.20,-3.95) {4};
  \node[case, noeudcase] at (2.00,-3.4) {10};
  \node[indice] at (2.00,-3.95) {5};
\end{tikzpicture}\\[2pt]
\small 
\end{minipage}
\end{minipage}
\caption{Le même tas vu comme un arbre et comme le tableau réellement stocké (indices en
gris). Le parent de la case 5 (10) est la case $(5-1)/2 = 2$ (20).}
\end{figure}

% ==============================================================================
\section{Insertion : on ajoute en bas, puis on remonte}
% ==============================================================================

\begin{lstlisting}[caption={Insertion (\texttt{inserer} et \texttt{monter}).}]
// Insertion : ajoute la valeur dans la premiere case libre (a la fin),
// puis la fait remonter. O(log n).
void inserer(int val) {
    tas.push_back(val);
    monter(tas.size() - 1);
}

// Fait remonter la case i : tant qu'elle est plus grande que son parent,
// on les echange. S'arrete a la racine ou quand le parent est plus grand.
void monter(int i) {
    while (i > 0 && tas[(i - 1) / 2] < tas[i]) {
        swap(tas[i], tas[(i - 1) / 2]);
        echanges++;
        i = (i - 1) / 2;
    }
}
\end{lstlisting}

La nouvelle valeur est placée dans la première case libre (à la fin du tableau). Tant
qu'elle est plus grande que son parent, on les \textbf{échange} : elle remonte d'un
niveau. Elle s'arrête quand son parent est plus grand, ou à la racine.

Exemple : on insère dans l'ordre $10, 4, 15, 20, 8, 30$. Voici l'insertion de \textbf{20}
dans le tas $[15, 4, 10]$ (en vert la valeur qui remonte, en orange la case comparée, en
rouge l'échange) :

\begin{figure}[H]
\centering
\begin{minipage}[t]{0.32\textwidth}\centering
\begin{tikzpicture}[scale=0.62, transform shape]
  \draw[arete] (0,0) -- (-1.6,-1.1);
  \draw[arete] (0,0) -- (1.6,-1.1);
  \draw[arete] (-1.6,-1.1) -- (-2.4,-2.2);
  \node[noeud] (n0) at (0,0) {15};
  \node[noeud] (n1) at (-1.6,-1.1) {4};
  \node[noeud] (n2) at (1.6,-1.1) {10};
  \node[neuf] (n3) at (-2.4,-2.2) {20};
  \node[case, noeudcase] at (-1.20,-3.4) {15};
  \node[indice] at (-1.20,-3.95) {0};
  \node[case, noeudcase] at (-0.40,-3.4) {4};
  \node[indice] at (-0.40,-3.95) {1};
  \node[case, noeudcase] at (0.40,-3.4) {10};
  \node[indice] at (0.40,-3.95) {2};
  \node[case, neufcase] at (1.20,-3.4) {20};
  \node[indice] at (1.20,-3.95) {3};
\end{tikzpicture}\\[2pt]
\small \textbf{1.} 20 est ajouté à la fin (case 3)
\end{minipage}
\hfill
\begin{minipage}[t]{0.32\textwidth}\centering
\begin{tikzpicture}[scale=0.62, transform shape]
  \draw[arete] (0,0) -- (-1.6,-1.1);
  \draw[arete] (0,0) -- (1.6,-1.1);
  \draw[arete] (-1.6,-1.1) -- (-2.4,-2.2);
  \node[noeud] (n0) at (0,0) {15};
  \node[actif] (n1) at (-1.6,-1.1) {4};
  \node[noeud] (n2) at (1.6,-1.1) {10};
  \node[neuf] (n3) at (-2.4,-2.2) {20};
  \draw[echange] (n3) to[bend right=35] (n1);
  \node[case, noeudcase] at (-1.20,-3.4) {15};
  \node[indice] at (-1.20,-3.95) {0};
  \node[case, actifcase] at (-0.40,-3.4) {4};
  \node[indice] at (-0.40,-3.95) {1};
  \node[case, noeudcase] at (0.40,-3.4) {10};
  \node[indice] at (0.40,-3.95) {2};
  \node[case, neufcase] at (1.20,-3.4) {20};
  \node[indice] at (1.20,-3.95) {3};
\end{tikzpicture}\\[2pt]
\small \textbf{2.} parent = case $(3-1)/2 = 1$ : $4 < 20$, on échange
\end{minipage}
\hfill
\begin{minipage}[t]{0.32\textwidth}\centering
\begin{tikzpicture}[scale=0.62, transform shape]
  \draw[arete] (0,0) -- (-1.6,-1.1);
  \draw[arete] (0,0) -- (1.6,-1.1);
  \draw[arete] (-1.6,-1.1) -- (-2.4,-2.2);
  \node[actif] (n0) at (0,0) {15};
  \node[neuf] (n1) at (-1.6,-1.1) {20};
  \node[noeud] (n2) at (1.6,-1.1) {10};
  \node[noeud] (n3) at (-2.4,-2.2) {4};
  \draw[echange] (n1) to[bend right=35] (n0);
  \node[case, actifcase] at (-1.20,-3.4) {15};
  \node[indice] at (-1.20,-3.95) {0};
  \node[case, neufcase] at (-0.40,-3.4) {20};
  \node[indice] at (-0.40,-3.95) {1};
  \node[case, noeudcase] at (0.40,-3.4) {10};
  \node[indice] at (0.40,-3.95) {2};
  \node[case, noeudcase] at (1.20,-3.4) {4};
  \node[indice] at (1.20,-3.95) {3};
\end{tikzpicture}\\[2pt]
\small \textbf{3.} parent = case 0 : $15 < 20$, on échange
\end{minipage}

\vspace{6mm}

\begin{minipage}[t]{0.32\textwidth}\centering
\begin{tikzpicture}[scale=0.62, transform shape]
  \draw[arete] (0,0) -- (-1.6,-1.1);
  \draw[arete] (0,0) -- (1.6,-1.1);
  \draw[arete] (-1.6,-1.1) -- (-2.4,-2.2);
  \node[neuf] (n0) at (0,0) {20};
  \node[noeud] (n1) at (-1.6,-1.1) {15};
  \node[noeud] (n2) at (1.6,-1.1) {10};
  \node[noeud] (n3) at (-2.4,-2.2) {4};
  \node[case, neufcase] at (-1.20,-3.4) {20};
  \node[indice] at (-1.20,-3.95) {0};
  \node[case, noeudcase] at (-0.40,-3.4) {15};
  \node[indice] at (-0.40,-3.95) {1};
  \node[case, noeudcase] at (0.40,-3.4) {10};
  \node[indice] at (0.40,-3.95) {2};
  \node[case, noeudcase] at (1.20,-3.4) {4};
  \node[indice] at (1.20,-3.95) {3};
\end{tikzpicture}\\[2pt]
\small \textbf{4.} 20 est à la racine : on s'arrête
\end{minipage}
\hfill
\begin{minipage}[t]{0.32\textwidth}\centering
\begin{tikzpicture}[scale=0.62, transform shape]
  \draw[arete] (0,0) -- (-1.6,-1.1);
  \draw[arete] (0,0) -- (1.6,-1.1);
  \draw[arete] (-1.6,-1.1) -- (-2.4,-2.2);
  \draw[arete] (-1.6,-1.1) -- (-0.8,-2.2);
  \node[noeud] (n0) at (0,0) {20};
  \node[actif] (n1) at (-1.6,-1.1) {15};
  \node[noeud] (n2) at (1.6,-1.1) {10};
  \node[noeud] (n3) at (-2.4,-2.2) {4};
  \node[neuf] (n4) at (-0.8,-2.2) {8};
  \node[case, noeudcase] at (-1.60,-3.4) {20};
  \node[indice] at (-1.60,-3.95) {0};
  \node[case, actifcase] at (-0.80,-3.4) {15};
  \node[indice] at (-0.80,-3.95) {1};
  \node[case, noeudcase] at (0.00,-3.4) {10};
  \node[indice] at (0.00,-3.95) {2};
  \node[case, noeudcase] at (0.80,-3.4) {4};
  \node[indice] at (0.80,-3.95) {3};
  \node[case, neufcase] at (1.60,-3.4) {8};
  \node[indice] at (1.60,-3.95) {4};
\end{tikzpicture}\\[2pt]
\small \textbf{Insérer 8} : parent 15 $> 8$, aucun échange
\end{minipage}
\hfill
\begin{minipage}[t]{0.32\textwidth}\centering
\begin{tikzpicture}[scale=0.62, transform shape]
  \draw[arete] (0,0) -- (-1.6,-1.1);
  \draw[arete] (0,0) -- (1.6,-1.1);
  \draw[arete] (-1.6,-1.1) -- (-2.4,-2.2);
  \draw[arete] (-1.6,-1.1) -- (-0.8,-2.2);
  \draw[arete] (1.6,-1.1) -- (0.8,-2.2);
  \node[neuf] (n0) at (0,0) {30};
  \node[noeud] (n1) at (-1.6,-1.1) {15};
  \node[noeud] (n2) at (1.6,-1.1) {20};
  \node[noeud] (n3) at (-2.4,-2.2) {4};
  \node[noeud] (n4) at (-0.8,-2.2) {8};
  \node[noeud] (n5) at (0.8,-2.2) {10};
  \node[case, neufcase] at (-2.00,-3.4) {30};
  \node[indice] at (-2.00,-3.95) {0};
  \node[case, noeudcase] at (-1.20,-3.4) {15};
  \node[indice] at (-1.20,-3.95) {1};
  \node[case, noeudcase] at (-0.40,-3.4) {20};
  \node[indice] at (-0.40,-3.95) {2};
  \node[case, noeudcase] at (0.40,-3.4) {4};
  \node[indice] at (0.40,-3.95) {3};
  \node[case, noeudcase] at (1.20,-3.4) {8};
  \node[indice] at (1.20,-3.95) {4};
  \node[case, noeudcase] at (2.00,-3.4) {10};
  \node[indice] at (2.00,-3.95) {5};
\end{tikzpicture}\\[2pt]
\small \textbf{Insérer 30} : 2 échanges (avec 10 puis 20)
\end{minipage}
\caption{Insertion pas à pas. Le tableau final est $[30, 15, 20, 4, 8, 10]$.}
\end{figure}

Une valeur remonte au plus de la hauteur de l'arbre, soit $\log_2 n$ niveaux (17 pour
200\,000 valeurs). Les $n$ insertions coûtent donc au pire $O(n \log n)$.

% ==============================================================================
\section{Extraction : on retire la racine, puis on redescend}
% ==============================================================================

\begin{lstlisting}[caption={Extraction du maximum (\texttt{extraireMax} et \texttt{descendre}).}]
// Extraction : retire et renvoie la racine (le max). Le dernier element
// prend sa place, puis on l'ajuste en le faisant descendre. O(log n).
int extraireMax() {
    int racine = tas[0];
    tas[0] = tas.back();
    tas.pop_back();
    if (!tas.empty()) descendre(0);
    return racine;
}

// Fait descendre la case i : on l'echange avec le plus grand de ses deux
// enfants tant qu'il est plus grand qu'elle. S'arrete quand elle est a sa
// place ou qu'elle n'a plus d'enfant.
void descendre(int i) {
    int n = tas.size();
    while (true) {
        int max = i;
        int g = 2 * i + 1;
        int d = 2 * i + 2;
        if (g < n && tas[g] > tas[max]) max = g;
        if (d < n && tas[d] > tas[max]) max = d;
        if (max == i) return;
        swap(tas[i], tas[max]);
        echanges++;
        i = max;
    }
}
\end{lstlisting}

Le maximum est la racine. Pour ne pas laisser de trou, on met le \textbf{dernier} élément
à la racine et on raccourcit le tableau. Cet élément est petit, donc la règle n'est plus
respectée : on l'\textbf{ajuste} en l'échangeant avec le plus grand de ses deux enfants,
jusqu'à ce qu'il soit plus grand que ses enfants ou qu'il n'en ait plus.

\begin{figure}[H]
\centering
\begin{minipage}[t]{0.32\textwidth}\centering
\begin{tikzpicture}[scale=0.62, transform shape]
  \draw[arete] (0,0) -- (-1.6,-1.1);
  \draw[arete] (0,0) -- (1.6,-1.1);
  \draw[arete] (-1.6,-1.1) -- (-2.4,-2.2);
  \draw[arete] (-1.6,-1.1) -- (-0.8,-2.2);
  \draw[arete] (1.6,-1.1) -- (0.8,-2.2);
  \node[actif] (n0) at (0,0) {30};
  \node[noeud] (n1) at (-1.6,-1.1) {15};
  \node[noeud] (n2) at (1.6,-1.1) {20};
  \node[noeud] (n3) at (-2.4,-2.2) {4};
  \node[noeud] (n4) at (-0.8,-2.2) {8};
  \node[neuf] (n5) at (0.8,-2.2) {10};
  \node[case, actifcase] at (-2.00,-3.4) {30};
  \node[indice] at (-2.00,-3.95) {0};
  \node[case, noeudcase] at (-1.20,-3.4) {15};
  \node[indice] at (-1.20,-3.95) {1};
  \node[case, noeudcase] at (-0.40,-3.4) {20};
  \node[indice] at (-0.40,-3.95) {2};
  \node[case, noeudcase] at (0.40,-3.4) {4};
  \node[indice] at (0.40,-3.95) {3};
  \node[case, noeudcase] at (1.20,-3.4) {8};
  \node[indice] at (1.20,-3.95) {4};
  \node[case, neufcase] at (2.00,-3.4) {10};
  \node[indice] at (2.00,-3.95) {5};
\end{tikzpicture}\\[2pt]
\small \textbf{1.} on retire la racine 30, le dernier (10) va prendre sa place
\end{minipage}
\hfill
\begin{minipage}[t]{0.32\textwidth}\centering
\begin{tikzpicture}[scale=0.62, transform shape]
  \draw[arete] (0,0) -- (-1.6,-1.1);
  \draw[arete] (0,0) -- (1.6,-1.1);
  \draw[arete] (-1.6,-1.1) -- (-2.4,-2.2);
  \draw[arete] (-1.6,-1.1) -- (-0.8,-2.2);
  \node[neuf] (n0) at (0,0) {10};
  \node[noeud] (n1) at (-1.6,-1.1) {15};
  \node[actif] (n2) at (1.6,-1.1) {20};
  \node[noeud] (n3) at (-2.4,-2.2) {4};
  \node[noeud] (n4) at (-0.8,-2.2) {8};
  \draw[echange] (n0) to[bend right=35] (n2);
  \node[case, neufcase] at (-1.60,-3.4) {10};
  \node[indice] at (-1.60,-3.95) {0};
  \node[case, noeudcase] at (-0.80,-3.4) {15};
  \node[indice] at (-0.80,-3.95) {1};
  \node[case, actifcase] at (0.00,-3.4) {20};
  \node[indice] at (0.00,-3.95) {2};
  \node[case, noeudcase] at (0.80,-3.4) {4};
  \node[indice] at (0.80,-3.95) {3};
  \node[case, noeudcase] at (1.60,-3.4) {8};
  \node[indice] at (1.60,-3.95) {4};
\end{tikzpicture}\\[2pt]
\small \textbf{2.} enfants 15 et 20 : le plus grand est 20 $> 10$, on échange
\end{minipage}
\hfill
\begin{minipage}[t]{0.32\textwidth}\centering
\begin{tikzpicture}[scale=0.62, transform shape]
  \draw[arete] (0,0) -- (-1.6,-1.1);
  \draw[arete] (0,0) -- (1.6,-1.1);
  \draw[arete] (-1.6,-1.1) -- (-2.4,-2.2);
  \draw[arete] (-1.6,-1.1) -- (-0.8,-2.2);
  \node[noeud] (n0) at (0,0) {20};
  \node[noeud] (n1) at (-1.6,-1.1) {15};
  \node[neuf] (n2) at (1.6,-1.1) {10};
  \node[noeud] (n3) at (-2.4,-2.2) {4};
  \node[noeud] (n4) at (-0.8,-2.2) {8};
  \node[case, noeudcase] at (-1.60,-3.4) {20};
  \node[indice] at (-1.60,-3.95) {0};
  \node[case, noeudcase] at (-0.80,-3.4) {15};
  \node[indice] at (-0.80,-3.95) {1};
  \node[case, neufcase] at (0.00,-3.4) {10};
  \node[indice] at (0.00,-3.95) {2};
  \node[case, noeudcase] at (0.80,-3.4) {4};
  \node[indice] at (0.80,-3.95) {3};
  \node[case, noeudcase] at (1.60,-3.4) {8};
  \node[indice] at (1.60,-3.95) {4};
\end{tikzpicture}\\[2pt]
\small \textbf{3.} 10 n'a plus d'enfant : on s'arrête, 20 est le nouveau max
\end{minipage}
\caption{Extraction du maximum du tas $[30, 15, 20, 4, 8, 10]$.}
\end{figure}

Une extraction descend au plus de $\log_2 n$ niveaux. Vider le tas demande $n$ extractions,
donc $O(n \log n)$. Les valeurs sortent dans l'ordre décroissant : c'est le principe du
\textbf{tri par tas}.

% ==============================================================================
\section{Résultats}
% ==============================================================================

Machine : Intel Core i5-1135G7, Linux, \texttt{g++} 15.2, \texttt{-std=c++17 -O3}.
Le fichier \texttt{data/data.txt} contient $3, 6, 9, \dots, 600\,000$ : 200\,000 valeurs
\textbf{triées par ordre croissant}. Le programme les utilise telles quelles, puis
mélangées (\texttt{std::shuffle}, graine 42). La lecture du fichier n'est pas chronométrée.

\begin{table}[H]
\centering
\begin{tabular}{l r r r r}
\toprule
& \multicolumn{2}{c}{\textbf{Temps}} & \multicolumn{2}{c}{\textbf{Échanges}} \\
\cmidrule(lr){2-3}\cmidrule(lr){4-5}
\textbf{Ordre des valeurs} & Insertions & Extractions & Insertions & Extractions \\
\midrule
Fichier (trié)  & 6,8 ms & 24,8 ms & 3\,137\,875 & 2\,937\,876 \\
Mélangées       & 6,2 ms & 32,5 ms &   256\,598 & 3\,001\,271 \\
\bottomrule
\end{tabular}
\caption{Sortie du programme (les temps varient d'environ 10\,\% d'une exécution à l'autre ;
les échanges sont toujours identiques).}
\end{table}

\textbf{Interprétation.}
\begin{itemize}
  \item \textbf{Insertions, fichier trié : le pire cas.} Chaque nouvelle valeur est plus
        grande que toutes les autres, donc elle remonte \textbf{jusqu'à la racine} :
        3\,137\,875 échanges, soit environ 15,7 par valeur (proche de $\log_2 200\,000 \approx 17$).
  \item \textbf{Insertions, valeurs mélangées.} Une valeur au hasard est en général plus
        petite que son parent et s'arrête tout de suite : 256\,598 échanges, environ
        1,3 par valeur, \textbf{12 fois moins}.
  \item \textbf{Pourtant le temps est presque le même} (6,8 contre 6,2 ms). Cela s'explique probablement ainsi : dans le cas trié,
        la boucle de \texttt{monter} se comporte toujours de la même façon (le processeur
        prédit parfaitement le test) et elle parcourt le haut du tableau, qui reste en
        cache. Dans le cas mélangé, la boucle s'arrête à un moment imprévisible.
        Avec 200\,000 valeurs, les deux phases d'insertion restent de toute façon très
        rapides.
  \item \textbf{Extractions : environ 3 millions d'échanges dans les deux cas.} Le dernier
        élément placé à la racine est petit et redescend presque jusqu'en bas, quel que
        soit l'ordre de départ : c'est bien $O(n \log n)$, environ 15 échanges par
        extraction. C'est la phase la plus lente (3,5 à 5 fois le temps des insertions), car chaque
        niveau demande deux comparaisons (les deux enfants) au lieu d'une.
\end{itemize}


% ==============================================================================
\section{Comparaison avec la STL}
% ==============================================================================

Le programme 11 fait les mêmes mesures (insérer les 200\,000 valeurs, puis retirer le
maximum jusqu'à vider la structure) avec trois conteneurs de la STL :

\begin{table}[H]
\centering
\small
\begin{tabular}{l l l l}
\toprule
\textbf{Structure} & \textbf{Organisation} & \textbf{Insérer} & \textbf{Retirer le max} \\
\midrule
Notre \texttt{TasMax}       & tas dans un \texttt{vector}     & \texttt{inserer(v)}  & \texttt{extraireMax()} \\
\texttt{priority\_queue}    & tas dans un \texttt{vector}     & \texttt{push(v)}     & \texttt{top()} + \texttt{pop()} \\
\texttt{set<int>}           & arbre rouge-noir, trié          & \texttt{insert(v)}   & \texttt{erase(prev(end()))} \\
\texttt{map<int,int>}       & arbre rouge-noir (valeur $\to$ nombre) & \texttt{m[v]++} & \texttt{prev(end())} : $-1$ ou effacer \\
\bottomrule
\end{tabular}
\caption{Les quatre structures. Toutes insèrent et retirent le max en $O(\log n)$.}
\end{table}

\texttt{set} ignore les doublons (le fichier n'en contient pas) ; \texttt{map} les compte,
comme un tas. Un premier passage non chronométré est fait avant les mesures : sans lui, la
première structure à nœuds mesurée payait 20 à 25~ms de plus (réservation de la mémoire
par le système).

\begin{table}[H]
\centering
\begin{tabular}{l r r r r}
\toprule
\textbf{Ordre + phase} & \textbf{Notre \texttt{TasMax}} & \textbf{\texttt{priority\_queue}} & \textbf{\texttt{set}} & \textbf{\texttt{map}} \\
\midrule
Fichier (trié), insertions   & 7,5 ms & 10 ms & 50 ms  & 50 ms \\
Fichier (trié), extractions  & 27 ms  & 21 ms & 12 ms  & 13 ms \\
\midrule
Mélangées, insertions        & 6,8 ms & 6 ms  & 105 ms & 140 ms \\
Mélangées, extractions       & 36 ms  & 28 ms & 35 ms  & 35 ms \\
\bottomrule
\end{tabular}
\caption{Valeurs typiques sur trois exécutions faites ensemble (programmes 10 et 11). Les
insertions mélangées de \texttt{set} et \texttt{map} varient de $\pm 20$\,\%.}
\end{table}

\begin{figure}[H]
\centering
\begin{tikzpicture}[x=0.55mm, y=5.5mm]
  \foreach \nom/\y/\ins/\ext/\tot in {
      Notre TasMax (trié)/7/7.5/27/35, Notre TasMax (mélangé)/6/6.8/36/43,
      priority\_queue (trié)/4.6/10/21/31, priority\_queue (mélangé)/3.6/6/28/34,
      set (trié)/2.2/50/12/62, set (mélangé)/1.2/105/35/140,
      map (trié)/-0.2/50/13/63, map (mélangé)/-1.2/140/35/175} {
    \node[anchor=east, font=\small] at (-2,\y) {\nom};
    \fill[chemin] (0,\y-0.35) rectangle (\ins,\y+0.35);
    \fill[orange] (\ins,\y-0.35) rectangle (\ins+\ext,\y+0.35);
    \node[anchor=west, font=\scriptsize] at (\ins+\ext,\y) {\tot\ ms};
  }
  \draw[thick, black!60] (0,-1.8) -- (0,7.6);
  \fill[chemin] (30,-2.7) rectangle (36,-2.3);
  \node[anchor=west, font=\small] at (38,-2.5) {Insertions};
  \fill[orange] (90,-2.7) rectangle (96,-2.3);
  \node[anchor=west, font=\small] at (98,-2.5) {Extractions};
\end{tikzpicture}
\caption{Temps total (insertions + extractions) pour 200\,000 valeurs.}
\end{figure}

\textbf{Interprétation.}
\begin{itemize}
  \item \textbf{\texttt{priority\_queue} est la même idée que notre tas} (un tas dans un
        \texttt{vector}) : les totaux sont proches (31--34~ms contre 35--43~ms). Son
        extraction est environ 20\,\% plus rapide, probablement parce que la version de la
        STL descend le trou jusqu'en bas en ne comparant que les deux enfants, puis remonte
        la valeur, au lieu de faire un échange complet à chaque niveau.
  \item \textbf{\texttt{set} et \texttt{map} insèrent 5 à 20 fois plus lentement.} Chaque
        insertion alloue un nœud (\texttt{new}), descend un arbre de pointeurs dispersés en
        mémoire, puis rééquilibre l'arbre. Le tas, lui, écrit dans un tableau contigu et,
        sur des valeurs mélangées, remonte à peine (1,3 échange en moyenne).
  \item \textbf{Leur extraction est rapide sur le fichier trié} (12--13~ms) : le max est
        le nœud le plus à droite, et comme les nœuds ont été alloués dans l'ordre croissant,
        ils se suivent en mémoire. Sur des valeurs mélangées les nœuds sont dispersés et
        l'extraction revient au niveau du tas (35~ms).
  \item \textbf{Bilan :} pour \og insérer puis toujours retirer le max \fg{}, un tas
        (\texttt{priority\_queue} ou le nôtre) est 2 à 5 fois plus rapide au total que
        \texttt{set}/\texttt{map}. Ces derniers ne valent la peine que si l'on a besoin
        de plus : parcourir les valeurs dans l'ordre, chercher ou supprimer une valeur
        quelconque en $O(\log n)$, ce qu'un tas ne sait pas faire.
\end{itemize}

% ==============================================================================
\section{Conclusion}
% ==============================================================================

Le tas max se range dans un simple tableau : le parent de la case $i$ est la case
$(i-1)/2$. L'insertion fait \textbf{remonter} la nouvelle valeur, l'extraction fait
\textbf{redescendre} le dernier élément placé à la racine, et chacune coûte au plus la
hauteur de l'arbre, $\log_2 n$. Sur 200\,000 valeurs, les insertions prennent environ 7~ms
et les extractions environ 25 à 33~ms. Des valeurs triées sont le pire cas des insertions
en nombre d'échanges (12 fois plus), sans gros écart de temps sur cette taille.
\texttt{std::priority\_queue} est la version STL de la même structure, un peu plus rapide
à l'extraction ; \texttt{std::set} et \texttt{std::map} donnent le même résultat mais sont
2 à 5 fois plus lents au total, à cause de l'allocation d'un nœud par valeur.

\end{document}
